% !PRINTED 70
\weberpage{70}
%<<CONT

If we call $\Theta$-functions with the same periods, of the same
order and the same characteristic, \gesperrt{related}, and if we further
designate a system of functions as \gesperrt{linearly dependent}
or \gesperrt{independent} according as a homogeneous linear relation
with constant (not all vanishing) coefficients
does or does not subsist between these functions, then the
theorem holds:

\gesperrt{$m + 1$ related $\Theta$-functions of order $m$ are
always linearly dependent;}

or:

\gesperrt{From $m$ linearly independent related $\Theta$-functions
 of order $m$ every other related
$\Theta$-function can be composed linearly (with constant
coefficients).}

This theorem is of the most fundamental importance for the
considerations that follow; from it there results, by \S\,18, 3.,
that every $T$-function of order $m$ can be composed, with the use of formula (3),
from $m$ linearly independent $\Theta$-functions.

By \S\,17, (11), the sum of the $m$ argument values for
which the function $\Theta_{g_1,\, g_2}(v)$ vanishes in the period parallelogram
is congruent, with respect to the modulus $(1, \omega)$, to
\[
m\,\frac{1+\omega}{2} + \frac{g_1\omega - g_2}{2}.
\]

The definition of the function $\vartheta$ given up to now leaves
undetermined a factor which may be a function of $\omega$. But
by an addition to the definition we can determine this factor
still more closely.

If the defining equations (1) are differentiated twice with respect to $v$
and once with respect to $\omega$, regarding $g_1, g_2$ as constant,
then a simple calculation shows that the function
\[
\frac{\partial^2 \vartheta(v,\omega)}{\partial v^2} - 4\pi i\,\frac{\partial \vartheta(v,\omega)}{\partial \omega}
\]
itself satisfies the conditions (1) and therefore has the form
\[
\varphi(\omega)\,\vartheta(v,\omega)
\]
in which $\varphi(\omega)$ is constant with respect to $v$, and hence a function of $\omega$
alone. If one now
\[
e^{-\frac{1}{4\pi i}\int \varphi(\omega)\,d\omega}\ \vartheta(v,\omega)
\]
%CONT>>
